\section{Problem, baseline, and the precise advance}
\label{sec:overview}

\subsection{What this continuation computes}

Let $\mathcal T$ be a finite triangulation of a compact orientable
three-manifold, and let $x\in\N^{7t}$ be an admissible normal vector, where
$t$ is the number of tetrahedra. Write $F(x)$ for the associated properly
embedded normal surface. Its coordinates are stored in binary. A single
coordinate can therefore describe more normal discs than could be explicitly
created in any feasible computation. Our output is the finite histogram
\[
 \Spec(x)(\chi,b,\epsilon)
 =\#\{C\in\pi_0F(x):\chi(C)=\chi,\ |\pi_0\partial C|=b,
                         \ C\text{ has orientability }\epsilon\},
\]
where $\epsilon$ distinguishes orientable and nonorientable components.
Multiplicity is an integer field, never an instruction to repeat a row.
For orientable components we also give $g=(2-\chi-b)/2$; for
nonorientable components we give the crosscap number $k=2-\chi-b$.

The mathematical constructions work in the stated orientable ambient setting
whenever the requisite normal-arc systems are available. The delivered geometry
validator has the narrower, inherited contract: a finite, compact, connected,
orientable triangulated three-manifold with exactly one torus boundary.
Generalized triangulations, repeated global vertices, and loop edges are allowed
when the validator's manifold checks succeed. Closed components of the supplied
surface are allowed. No irreducibility hypothesis is imposed.

The distinction between \emph{type} and \emph{embedding} is essential.
Abstract surface type does not determine an embedding or its peripheral data.
The output does not assign an individual normal vector to each component,
provide cutting instructions, or identify a disc as a compressing disc. In
particular, a vertex-link disc and a meridian disc both have $(\chi,b)=(1,1)$.
The existing componentwise boundary-homology query remains the appropriate
operation for essential-disc counts on the supported torus boundary.

\subsection{Repository state and reproducibility}

The baseline is the public ProveIt repository at commit
\begin{center}\small\baseline.\end{center}
All source retrievals used this explicit ref after resolving the current main
branch. The relevant subtree is
\path{Topology/UnknotRecognition}. Its maintained \path{fast/fastunknot}
package already contains exact recognition, compressed interval orbit counts,
normal-surface validation, componentwise weighted queries, and several
independent certificate checkers. The new work is additive: it introduces six
runtime modules and three focused test modules without altering the existing
recognition pipeline or its default choices. The accompanying manifest records
the exact baseline and delivered-file hashes.

Three earlier developments are particularly relevant. The maintained
\path{synthesis/normal_multiplicity.tex} derives aggregate orientability and
component-count formulas under common normal scaling. The maintained
\path{synthesis/weighted_components.tex}, incorporating report 53, gives
componentwise Euler/boundary-homology weights and a quadrilateral core for
essential-disc counts. The \code{coordinates} mode of
\code{normal\_component\_census} returns component normal vectors in dimension
$7t$. Its cheaper \code{summary} mode records boundary \emph{intersection
points}, not the number of boundary circles of each component
\cite{proveit-baseline}.

The new interface exposes the joint distribution of component topology, beyond
the separate totals of Euler characteristic, boundary components, or
orientability. Separate totals cannot determine this joint distribution.
For example, two tori and a sphere disjoint from a closed genus-two surface
both have two orientable components, Euler total zero, and no boundary, but
their type spectra differ. The baseline can already recover the same abstract
spectrum by composing full-coordinate recovery with a topology query for each
distinct component coordinate row. Our benchmark retains this exact composed
route as a reference.

\subsection{Relationship to classical results}

Agol--Hass--Thurston give polynomial encoded orbit counting, weighted orbit
counting, and component-topology computation for normal surfaces. In their
Corollary 17, the weights recover all $7t$ normal coordinates of a component;
their orientability test uses doubled normal coordinates in an orientable
ambient manifold \cite{aht}. Accordingly, neither the existence of a polynomial
topology algorithm nor the normal-double orientability principle is claimed
as new here. Our explicit refinement replaces coordinate recovery by a
two-coordinate observer, proves a small interval representation for its
boundary markers, inverts the resulting two histograms, and extends the
existing numerical core to all abstract component types. Priority for this
specific combination is not asserted; the primary-source comparison establishes
the nearest classical baseline.

The dimensions should not be mistaken for a universal speedup theorem.
The two-coordinate approach requires a boundary transversal and two weighted
surface queries, and the transversal can introduce more initial weight runs.
An input on which coordinate recovery is already cheap may therefore favour
the older route. The operation counts, the core-sensitive bound, and the paired
measurements below make these tradeoffs explicit.

\subsection{The global quasi-polynomial objective}

The project's intended objective is a practical exact implementation with a
quasi-polynomial worst-case recognition bound. We retain that goal without
changing the meaning of the complexity variable: $n$ denotes crossing number
for a knot-diagram algorithm, whereas $t$ and $B$ below describe a supplied
triangulation and binary normal vector.

There is also a quantitative distinction in the literature. Lackenby's 2021
Oxford slides state a recognition bound
\[
 2^{O((\log n)^3)}=n^{O((\log n)^2)}
\]
and describe a hierarchy depth of order $(\log n)^2$ \cite{lackenby-talk}.
This belongs to the general quasi-polynomial class
$2^{(\log n)^{O(1)}}$; it is weaker than the narrower target
$n^{O(\log n)}=2^{O((\log n)^2)}$. His July 2026 preprint develops hierarchy
algorithms but explicitly leaves some local operations insufficiently specified
for a running-time estimate; Section 9 bounds loop counts in terms of hierarchy
length and pattern complexity \cite{lackenby2026}. These sources do not supply
a completed costed implementation of the announced accelerated hierarchy.

The 2026 preprint also proves polynomial encoding and verification of
essential-hierarchy certificates under its stated hypotheses (Theorem 14.1).
That theorem is an important model for implementation; it does not give a
quasi-polynomial deterministic discovery procedure for those certificates.
The implementation tasks below should therefore build on its encoding
result, with discovery and total execution cost tracked separately
\cite{lackenby2026}.

Our new census has a polynomial encoded bound for one supplied vector.
This prevents an unnecessary expansion bottleneck inside a future hierarchy,
but it cannot bound the number of candidate vectors visited or the total
encoded size of successive cut manifolds. No complete recognizer in this
package acquires a new general asymptotic bound.

\subsection{Results at a glance}

\begin{center}
\begin{tabular}{p{0.29\linewidth}p{0.61\linewidth}}
\toprule
Operation & Proven and implemented result\\
\midrule
Original orbit representatives & Least points of all orbits in at most $G$
intervals, for $G$ static gap intervals in a complete trace.\\
Boundary distribution & A marker weight counts actual boundary circles inside
each surface component.\\
Full type spectrum & One boundary discovery and two weighted surface discoveries;
weight dimension two; disconnected and one-sided cases included.\\
Histogram inversion & At fixed weight dimension, ordered-map recovery uses
$O(s\log(s+1))$ comparisons and $O(s)$ integer operations; zero has a
separate equation.\\
Coordinate core & Peel vertex links, divide by quadrilateral content, query the
small core, and restore the full spectrum with parity-aware scaling.\\
Evidence & Source-bound producer certificates, independently implemented replay,
bounded Regina comparisons, mutations, and paired measurements.\\
\bottomrule
\end{tabular}
\end{center}

\section{Encoded normal surfaces and inherited primitives}
\label{sec:model}

\subsection{Normal coordinates and the size model}

Each tetrahedron has four triangle types and three quadrilateral types. A
normal vector is admissible if all entries are nonnegative integers, matching
equations hold across glued faces, and at most one quadrilateral type is
nonzero in each tetrahedron. The vector may be zero. We measure its numerical
size by
\[
 B=\max\bigl(1,\max_i\bits(x_i)\bigr).
\]
The combinatorial triangulation encoding has size $O(t\log(t+1))$, and the
coordinate encoding has size $O(tB)$. Let $N$ be the number of intersections
with global triangulation edges. Since every normal disc has three or four
corners, $N\le4\sum_i x_i$, and
\[
 \log(N+1)=O(B+\log(t+1)).
\]
None of the algorithms below has a loop bounded by $N$ or by $\sum_i x_i$.

Let $\M(b)$ denote an upper bound for multiplying two $b$-bit integers.
Additions, comparisons, division, gcd, histogram keys, and output multiplicities
are charged at their binary cost. In particular, eliminating a large common
factor from an orbit query does not eliminate the cost of reading that factor
in the original input.

\subsection{Interval pairings}

An interval pairing identifies equal-length integer intervals by a translation
or reflection. The mathematical universe is $[0,N)\cap\Z$. The inherited
\code{IntervalPairing} stores its domain and image endpoints inclusively;
all representative and weight intervals in this extension are half open.
Thus a stored source interval $[a,b]$ becomes the weight interval $[a,b+1)$.
The implementation and its tests keep this conversion explicit.

The normal arc construction gives $O(t)$ interval pairings on the global
edge-point universe. A connected orbit is exactly a connected surface
component. Applying the same construction only to boundary faces and boundary
edges gives another interval system whose orbits are boundary circles.
The order on both universes is fixed by the same global edge orientations.
This matters for one-vertex and loop-edge triangulations: vertex labels alone
do not determine the order along an edge.

We use the maintained AHT discovery engine and its independently replayable
local trace. Its operations are identity deletion, reflection trimming,
periodic merging, transmission, contraction of static intervals, and suffix
truncation. We use only their established local correctness and polynomial
encoded bound. For clarity, a static point has no remaining nontrivial pairing
incident to it and is a singleton of the \emph{current} relation; it may stand
for many original points removed earlier.

\subsection{Weighted histograms}

For an additive, piecewise-constant integer-vector weight $w$ on the point
universe, weighted AHT returns
\[
 H(v)=\#\left\{O:\sum_{p\in O}w(p)=v\right\}.
\]
Overlapping input intervals add their weights. Output keys are unique vectors
with positive integer multiplicities. Weight transport follows the local trace:
deleted points transfer their weights to retained equivalent points, and a
static contraction emits complete orbit weights. The existing independent
checker verifies the exact relation, weight source, transport, and histogram.
The extension uses this API directly; it does not substitute mere conservation
of total weight for orbitwise correctness.

For proofs involving a particular trace, let $E$ be its number of local events
and let $G$ be the sum, over contractions, of the number of static gap intervals.
Both count encoded records, not the numbers of represented points. With $k$
initial pairings, the maintained reductions do not increase the pairing count,
so each contraction has at most $2k+1$ static gaps; in particular
$G\le(2k+1)E$. Polynomial trace length makes all the bounds involving $E,G$
below polynomial in encoded input size.
